\subsection{伴随}

命题 1. --- 设 E 与 F 为两个希尔伯特空间。对每个映射 \(u \in \mathcal{L}(E; F)\)，存在唯一的映射 \(u^{*} \in \mathcal{L}(F; E)\)，使得

\[
\langle u (x) | y \rangle_ {\mathrm{F}} = \langle x | u ^ {*} (y) \rangle_ {\mathrm{E}}\tag{1}
\]

对所有 \(x \in E\) 与所有 \(y \in F\) 成立。映射 \(u \mapsto u^{*}\)（从 \(\mathcal{L}(\mathrm{E}, \mathrm{F})\) 到 \(\mathcal{L}(\mathrm{F}; \mathrm{E})\) 中）是双射、等距且半线性的（关于 K 的自同构 \(\xi \mapsto \bar{\xi}\)）。

设 \(\mathcal{S}(\mathrm{E},\mathrm{F})\) 为 \(\mathbf{E}\times \mathbf{F}\) 上所有连续半双线性型组成的空间，赋以范数

\[
\| \Phi \| = \sup _ {\| x \| \leqslant 1, \| y \| \leqslant 1} | \Phi (x, y) |.\tag{2}
\]

我们类似地定义空间 \(\mathcal{S}(\mathrm{F},\mathrm{E})\)。我们曾定义（V, p.~16, 推论 2）一个从 \(\mathcal{L}(\mathrm{E};\mathrm{F})\) 到 \(\mathcal{S}(\mathrm{F},\mathrm{E})\) 上的 Banach 空间同构，记作 \(u\mapsto\Phi_{u}\)，它由下式刻画

\[
\Phi_ {u} (y, x) = \langle y | u (x) \rangle_ {\mathrm{F}} (x \in \mathrm{E}, y \in \mathrm{F}).\tag{3}
\]

以类似的方式我们定义一个从 \(\mathcal{L}(\mathrm{F},\mathrm{E})\) 到 \(\mathcal{S}(\mathrm{E},\mathrm{F})\) 上的同构。最后，我们用下式定义映射 \(\Phi \mapsto \Phi^{*}\)（从 \(\mathcal{S}(\mathrm{F},\mathrm{E})\) 到 \(\mathcal{S}(\mathrm{E},\mathrm{F})\) 上）

\[
\Phi^ {*} (x, y) = \overline {{\Phi (y , x)}} \quad (x \in \mathrm{E}, y \in \mathrm{F}).\tag{4}
\]

这一映射是双射、半线性且等距的。但公式 (1) 可以改写为 \(\Phi_{u^{*}} = (\Phi_{u})^{*}\)，由此即得命题。

定义 1. --- 设 E 与 F 为两个希尔伯特空间。对每个连续线性映射 \(u: \mathbf{E} \to \mathbf{F}\)，由公式 (1) 定义的从 F 到 E 中的连续线性映射称为 \(u\) 的伴随，记作 \(u^*\)。

我们有

\[
(u + v) ^ {*} = u ^ {*} + v ^ {*}\tag{5}
\]

\[
(\lambda u) ^ {*} = \overline {{{{\lambda}}}} u ^ {*}\tag{6}
\]

\[
(u ^ {*}) ^ {*} = u\tag{7}
\]

\[
(1 _ {\mathrm{E}}) ^ {*} = 1 _ {\mathrm{E}}\tag{8}
\]

\[
(w u) ^ {*} = u ^ {*} w ^ {*};\tag{9}
\]

在所有这些公式中，u 与 v 属于 \(\mathcal{L}(\mathrm{E};\mathrm{F})\)，\(\lambda\) 属于 K，w 属于 \(\mathcal{L}(\mathrm{F};\mathrm{G})\)，其中 G 是一个希尔伯特空间。公式 (5) 与 (6) 表明 \(u\mapsto u^{*}\) 是半线性的。公式 (8) 是显然的。为证 (7)，我们对 (1) 的两边取共轭，得 \(\langle u^{*}(y)|x\rangle = \langle y|u(x)\rangle\)，这就证明 u 是 \(u^{*}\) 的伴随。最后，采用 (9) 的记号，对所有 \(z \in G\) 我们有

\[
\langle w (u (x)) | z \rangle = \langle u (x) | w ^ {*} (z) \rangle = \langle x | u ^ {*} (w ^ {*} (z)) \rangle ,
\]

因此 \(u^{*}w^{*}\) 是 wu 的伴随。

设 \(u: \mathrm{E} \to \mathrm{F}\) 为双射且连续的线性映射；于是它也是双连续的（I, p.~19, 推论 1）。由 (8) 与 (9) 我们立即推出 \(u^*\) 是双射且双连续的，并且

\[
(u ^ {- 1}) ^ {*} = (u ^ {*}) ^ {- 1}.\tag{10}
\]

命题 2. --- 对每个 \(u \in \mathcal{L}(\mathrm{E};\mathrm{F})\)，我们有

\[
\| u ^ {*} u \| = \| u u ^ {*} \| = \| u \| ^ {2} = \| u ^ {*} \| ^ {2}.\tag{11}
\]

由命题 1，\(\| u^{*}\| = \| u\|\)，因此 \(\| u^{*}u\| \leqslant \| u^{*}\| \cdot \| u\| \leqslant \| u\|^2\)。另一方面，

\[
\| u \| ^ {2} = \sup _ {\| x \| \leqslant 1} \| u (x) \| ^ {2} = \sup _ {\| x \| \leqslant 1} \langle u (x) | u (x) \rangle = \sup _ {\| x \| \leqslant 1} \langle x | u ^ {*} u (x) \rangle \leqslant \| u ^ {*} u \|,
\]

因此 \(\| u^{*}u\| = \| u\|^2\)。把 \(u\) 换成 \(u^{*}\)，得 \(\| uu^{*}\| = \| u^{*}\|^2\)，再由 \(\| u\| = \| u^{*}\|\) 即得 (11)。

设 \(E_{1}, \ldots, E_{n}\) 与 \(F_{1}, \ldots, F_{n}\) 为希尔伯特空间，且对 1 与 n 之间的每个整数 i，设 \(u_{i}\) 为从 \(E_{i}\) 到 \(F_{i}\) 中的连续线性映射。于是

\[
(u _ {1} \hat {\otimes} _ {2} \dots \hat {\otimes} _ {2} u _ {n}) ^ {*} = u _ {1} ^ {*} \hat {\otimes} _ {2} \dots \hat {\otimes} _ {2} u _ {n} ^ {*}.\tag{12}
\]

设 \(v\) 为连续线性映射 \(u_1 \hat{\otimes}_2 \ldots \hat{\otimes}_2 u_n\)，它从

\[
\mathbf {E} = \mathbf {E} _ {1} \hat {\otimes} _ {2} \dots \hat {\otimes} _ {2} \mathbf {E} _ {n} \quad \text {into} \quad \mathbf {F} = \mathbf {F} _ {1} \hat {\otimes} _ {2} \dots \hat {\otimes} _ {2} \mathbf {F} _ {n}
\]

而 w 为从 F 到 E 中的连续线性映射 \(u_{1}^{*}\hat{\otimes}_{2}\ldots\hat{\otimes}_{2}u_{n}^{*}\)。只需证明等式 \(\langle y|v(x)\rangle=\langle w(y)|x\rangle\)（对 \(x\in E\) 与 \(y\in F\)）。由线性性与连续性，可归结为 x 与 y 具有下列形式的情形

\[
x = x _ {1} \otimes \ldots \otimes x _ {n}, y = y _ {1} \otimes \ldots \otimes y _ {n}
\]

其中 \(x_{i} \in E_{i}\) 与 \(y_{i} \in F_{i}\)（对 \(1 \leqslant i \leqslant n\)）。于是由张量积中内积的定义（V, p.~27, 公式 (6)），得

\[
\langle y | v (x) \rangle = \prod_ {i = 1} ^ {n} \left\langle y _ {i} \mid u _ {i} \left(x _ {i}\right) \right\rangle = \prod_ {i = 1} ^ {n} \left\langle u _ {i} ^ {*} \left(y _ {i}\right) \mid x _ {i} \right\rangle = \left\langle w (y) \mid x \right\rangle .
\]

这就证明了我们的断言。

设 E 与 F 为两个希尔伯特空间，\(u \in \mathcal{L}(\mathrm{E}; \mathrm{F})\)，n 为正整数。若在公式 (12) 中令 \(u_{1} = \cdots = u_{n} = u\)，我们就得到如下结果：连续线性映射 \(\hat{\mathbf{T}}^{n}(u^{*})\)（从 \(\hat{\mathbf{T}}^{n}(\mathrm{F})\) 到 \(\hat{\mathbf{T}}^{n}(\mathrm{E})\) 中）是连续线性映射 \(\hat{\mathbf{T}}^{n}(u)\)（从 \(\hat{\mathbf{T}}^{n}(\mathrm{E})\) 到 \(\hat{\mathbf{T}}^{n}(\mathrm{F})\) 中）的伴随。公式

\[
\hat {\mathbf {S}} ^ {n} (u) ^ {*} = \hat {\mathbf {S}} ^ {n} (u ^ {*}), \quad \hat {\symbf {\Lambda}} ^ {n} (u) ^ {*} = \hat {\symbf {\Lambda}} ^ {n} (u ^ {*})\tag{13}
\]

可以仿照公式 (12) 来证明，只需用到 \(\hat{\mathbf{S}}^n (\mathrm{E})\)（V, p.~30, 公式 (15)）与 \(\hat{\Lambda}^n (\mathrm{E})\)（V, p.~33, 公式 (26)）中内积的定义。

注 1. --- 假定希尔伯特空间 E 不化为 0。我们把 \(\mathcal{L}(\mathbf{K};\mathbf{E})\) 与 E（借助映射 \(u\mapsto u(1)\)）等同；换言之，E 中的向量 \(x\) 与从 K 到 E 中的映射 \(\lambda\mapsto\lambda.x\) 等同。这时 \(x\) 的伴随是映射 \(x^{*}:E\to K\)（由 \(x^{*}(y)=\langle x|y\rangle\) 给出）。换句话说，\(x\mapsto x^{*}\) 是从 E 到其对偶上的典范半线性映射（V, p.~15）。

类似地，我们把数 \(\lambda \in \mathbf{K}\) 与 E 的自同态 \(\lambda. 1_{\mathrm{E}}\) 等同。这时 \(\lambda^{*}\) 恰好就是 \(\lambda\) 的共轭。

借助这些等同，我们可以定义乘积 \(t_{1} \ldots t_{n}\)，其中每个 \(t_{i}\) 或者是 K 中的数，或者是 E 中的向量，或者是属于 \(E'\) 的线性型，或者是 \(\mathcal{L}(E)\) 的元素，只要相邻的两个因子 \(t_{i}\) 与 \(t_{i+1}\) 从不属于下列类型之一：

\(\bullet\) \(xy\)，其中 \(x, y\) 都在 \(\mathbf{E}\) 中，或都在 \(\mathbf{E}'\) 中；

\(\bullet\) \(x\mathbf{A}\) 或 \(\mathbf{A}x'\)，其中 \(\mathbf{A} \in \mathcal{L}(\mathbf{E})\)，\(x \in \mathbf{E}\)，\(x' \in \mathbf{E}'\)。

我们有下列合成规则：

\begin{enumerate}
\def\labelenumi{\alph{enumi})}
\item
  结合性；
\item
  \(\mathbf{K}\) 的每个元素与其余所有因子交换；
\item
  我们有 \((t_1 \ldots t_n)^* = t_n^* \ldots t_1^*\)；换言之，乘积的伴随是诸伴随按相反顺序的乘积。此外 \(t^{**} = t\)。
\end{enumerate}

例如，设 \(x, y\) 属于 E，A 属于 \(\mathcal{L}(\mathrm{E})\)。于是 \(x^{*}y\) 表示内积 \(\langle x|y\rangle\)，而 \(x^{*}\mathrm{A}y\) 表示内积 \(\langle x|\mathrm{A}y\rangle\)。我们还有 \((\mathrm{A}^{*}x)^{*} = x^{*}\mathrm{A}^{**} = x^{*}\mathrm{A}\)，因此 \((\mathrm{A}^{*}x)^{*}y = x^{*}\mathrm{A}y\)，它可以解释为

\[
\langle \mathrm{A} ^ {*} x | y \rangle = \langle x | \mathrm{Ay} \rangle
\]

这与伴随的定义一致。我们注意到 \(yx^{*}\) 是 E 的自同态 \(z \mapsto y\langle x|z\rangle\)，因为由结合性，\(yx^{*}z\) 既可以解释为 \(y(x^{*}z)\)，也可以解释为 \(y.\langle x|z\rangle\)。

依照 Dirac \(^{1}\)，在数学物理的大多数著作中，E 的元素用符号 \(|x\rangle\) 表示，E' 的元素用 \(\langle t|\) 表示。内积写作 \(\langle x|y\rangle = \langle x|.|y\rangle\)，而乘积中的第一条禁用规则排除了记号 \(\rangle|\) 与 \(|\langle\) 的组合，例如 \(|x\rangle|y\rangle\)。

\(^{1}\) 参见 P. A. M. DIRAC, \emph{Quantum Mechanics}, Oxford University Press, New York, 1935.

命题 3. --- 设 E 与 F 为两个希尔伯特空间，\(u \in \mathcal{L}(\mathrm{E};\mathrm{F})\)。下列条件等价：

\begin{enumerate}
\def\labelenumi{(\roman{enumi})}
\item
  \(u\) 是拓扑向量空间同构，且其逆等于 \(u^{*}\)；
\item
  \(u\) 是满射且 \(u^{*}u = 1_{\mathrm{E}}\)；
\item
  \(u\) 是单射且 \(uu^{*} = 1_{\mathrm{F}}\)；
\item
  \(u\) 是赋范空间同构；
\item
  \(u\) 是希尔伯特空间同构。
\end{enumerate}

条件 (1) 意味着我们既有 \(u^{*}u = 1_{\mathrm{E}}\) 又有 \(uu^{*} = 1_{\mathrm{F}}\)。于是 (i)、(ii) 与 (iii) 的等价性由 E, II, § 3, No.~8, 命题 8 得出。(iv) 与 (v) 的等价性我们已经在（V, p.~5）中见过。最后，关系 \(u^{*}u = 1_{\mathrm{E}}\) 等价于 \(\langle x|u^{*}u(y)\rangle = \langle x|y\rangle\)，即 \(\langle u(x)|u(y)\rangle = \langle x|y\rangle\)（对 E 中所有 \(x,y\)），并且显然蕴含 \(u\) 是单射；这就证明了 (ii) 与 (v) 的等价性。

希尔伯特空间 E 的自同构也称为酉算子，即算子 \(u \in \mathcal{L}(\mathrm{E})\)（满足 \(uu^{*} = u^{*}u = 1_{E}\)）。

注 2. --- 关系 \(u^{*}u = 1_{\mathrm{E}}\) 并不刻画希尔伯特空间 E 的所有自同构。例如，设 \(\mathrm{E} = \ell^2 (\mathbf{N})\)，并设 \(u\) 由 \(u(x_{n}) = x_{n - 1}\)（\(n\geqslant 1\)）及 \(u(x)_0 = 0\) 定义。我们有 \(\| u(x)\| = \| x\|\)（对所有 \(x\in \mathrm{E}\)），即 \(u^{*}u = 1_{\mathrm{E}}\)，但 \(u\) 不是满射。

注 3. --- 伴随 \(u^{*}\) 的定义 (1) 也可以写成

\[
\langle y | u (x) \rangle = \langle u ^ {*} (y) | x \rangle ,
\]

或者，依 V, p.~15，写成

\[
\langle u (x), y ^ {*} \rangle = \langle x, (u ^ {*} (y)) ^ {*} \rangle .
\]

但我们也有 \(\langle u(x), y^* \rangle = \langle x, {}^t u(y^*) \rangle\)，因此我们可以用转置来表示伴随：

\[
(u ^ {*} (y)) ^ {*} = ^ {t} u (y ^ {*}).
\]

